\scanpage{5}
\mspage{2}
\margin{\diagram{\begin{tikzcd}[ampersand replacement=\&,column sep=small,row sep=small]
K\arrow[r,dash]\arrow[dr,dash]\&K'\arrow[dr,dash]\&{}\\
{}\&\bar K\arrow[r,dash]\&\bar K'
\end{tikzcd}}}
$K_1$ de $\bar K/K$, clôture \inserted{alg.} sép. de $K$ dans $K'$, i.e. $K_1=\bar K\cap K'$.

Quand $K'$ est une extension de type fini de $K$, $K_1$ est une ext. finie de $K$, donc on conclut que l'image de (3) est \unsure{tjs} un sous-groupe d'indice fini, égale à $E_{\bar K/K}$ ssi \inserted{$K_1=K$ i.e.} $K$ est séparablement alg. clos dans $K'$. D'ailleurs, on montre sans mal que \inserted{(si $K'$ est ext. de type fini)} l'hom (3) est injectif ssi $K'$ est une extension algébrique de $K$. Donc il est bijectif ssi $K'$ est une extension radicielle de $K$. Dans la suite nous nous bornerons (provisoirement) aux corps de car. $0$, et la condition précédente \uline{signifie} \inserted{donc} que $K\unsure{-|\longrightarrow}K'$ est un \uline{isomorphisme}.

Ainsi, le foncteur $K\longrightarrow B_K$ ou $K\unsure{-|\longrightarrow}\Pi_K$, ou $(K,\bar K)\unsure{-|\longrightarrow}E_{\bar K/K}$, est \uline{conservatif} quand on se limite comme morphismes de corps \inserted{$K\longrightarrow K'$ (de car. $0$)} à ceux \deleted{pour} qui font de $K'$ une ext. de type fini de $K$.
\margin{4}
